% Generated draft: requires mathematical and source review.
\documentclass[11pt]{article}
\usepackage[margin=25mm]{geometry}
\usepackage{fontspec}
\setmainfont{Latin Modern Roman}
\usepackage{amsmath,amssymb,mathtools,unicode-math}
\setmathfont{Latin Modern Math}
\usepackage{enumitem,needspace,xurl,hyperref}
\setlength{\parindent}{0pt}
\setlength{\parskip}{5pt}
\setlength{\emergencystretch}{3em}
\newcommand{\N}{\mathbb{N}}
\newcommand{\Z}{\mathbb{Z}}
\newcommand{\Q}{\mathbb{Q}}
\newcommand{\R}{\mathbb{R}}
\newcommand{\C}{\mathbb{C}}
\newcommand{\F}{\mathbb{F}}
\newcommand{\A}{\mathbb{A}}
\newcommand{\PP}{\mathbb{P}}
\newcommand{\E}{\mathbb{E}}
\newcommand{\eps}{\varepsilon}
\newcommand{\dd}{\,\mathrm{d}}
\begin{document}
\section*{UnsolvedMath ID: 30001849}
\textbf{Source problem code:} catalogue entry
\subsection{Definitions and mathematical statement}
Let \(p\) be a prime. A finite \(p\)-group is a finite group whose order is a power of \(p\). For a finite \(p\)-group \(H\), its generator rank is
\[
d(H)=\dim_{\F_p}H/\bigl(H^p[H,H]\bigr),
\]
where \(H^p\) is the subgroup generated by the \(p\)th powers and \([H,H]\) is the commutator subgroup. Write \(\operatorname{Aut}(H)\) for its automorphism group.

The supplied conjecture concerns a finite \(p\)-group \(G\) that is either a Schur \(\sigma\)-group of rank \(g\), or a maximal \(p\)-class-\(c\) quotient of such a group. It considers imaginary quadratic fields \(K\) for which \(A_K\) has rank \(g\). Retaining the source notation, put
\[
M(G;g,h,p)=\frac{z(G)^g}{|\operatorname{Aut}(G)|}\,p^{-gh}
\prod_{k=1}^{g}\bigl(p^g-p^{g-k}\bigr)
\prod_{k=1}^{h}\bigl(p^g-p^{h-k}\bigr).
\]
The source asks whether the probability that \(G_K\cong G\), in the Schur \(\sigma\)-group case, equals \(M(G;g,h,p)\), and whether the probability that the maximal \(p\)-class-\(c\) quotient of \(G_K\) is isomorphic to \(G\), in the quotient case, equals the same expression. Both probabilities are conditioned on \(A_K\) having rank \(g\). The source also describes the proposed probability as a limiting frequency.

\emph{Unresolved qualifications.} This is a source-level conjectural template, not a fully specified limiting-frequency assertion. The supplied record does not define \(A_K\), \(G_K\), \(z(G)\), or \(h\); specify the intended Schur \(\sigma\)-group convention; or define the filtration used for maximal \(p\)-class-\(c\) quotients. The generator-rank definition above supplies standard terminology, but the record does not explicitly identify it as its rank convention. No restriction to odd primes, admissible range of \(c\) or \(h\), ordering of imaginary quadratic fields, or probability measure is supplied. These choices must be recovered before the displayed expression determines a precise conjecture. In particular, no definitions of \(z(G)\) or \(h\), and no field-ordering convention, are inferred here.
\subsection{Short English statement}
The supplied conjecture predicts a particular frequency for a finite Schur \(\sigma\)-group, or a specified quotient of one, among groups attached to imaginary quadratic fields with fixed class-group rank in the source notation. Several essential symbols, group conventions, and the meaning of limiting frequency are unspecified, so the formula is preserved with these gaps explicit. The record labels the problem open; this draft does not independently verify that status.
\subsection{Sources}
\begin{itemize}
\item Supplied local UnsolvedMath registry record, permanent ID \(30001849\), exact source problem code \texttt{catalogue entry}, title: ``Schur \(\sigma\)-Group Frequency Heuristics.'' The record gives proposed year \(2011\) and no named proposer.
\item Supplied report citation: \emph{Explicit Methods in Number Theory (2012)}, Oberwolfach Reports, pp.\ 1971--2044, DOI \texttt{10.4171/owr/2011/35}: \url{https://doi.org/10.4171/owr/2011/35}. The record describes Oberwolfach Reports as documenting mathematical workshops and their open-problem discussions. The supplied citation's year \(2012\) and the record's proposed year \(2011\) are retained without reconciliation.
\item Supplied literature source, labelled ``Boston--Bush--Hajir heuristics'': \url{https://doi.org/10.1007/s00208-016-1449-3}.
\item Supplied literature source, labelled ``original OWR report'': \url{https://doi.org/10.4171/owr/2011/35}.
\item Literature assessment dated exactly \texttt{2026-08-21}: the supplied curation labels the problem open and describes it as an isolated conjectural fragment of the non-abelian Cohen--Lenstra heuristic, with no theorem establishing the stated frequency law in general. This is supplied dated curation, not an independently verified literature assessment. The background expressly requests independent verification of time-sensitive status and cited claims.
\item The registry assigns the statement status \texttt{corrected\_verified} and reports that its mathematical content was checked against the official Oberwolfach report, with notation, context, and question form normalized without changing the source-level claim. This is a supplied verification claim; no source visit or independent verification is asserted here.
\end{itemize}
\end{document}
